% Unidad U014. Biografía estructural completa del carácter de autoescala.

\label{gfp:b:chap:biografia-completa-phi}

\section{Finite Selection of the Self-Scaling Seed}

Let \(\mathcal U_6\subset\{0,\ldots,999\}^6\) be the catalog of the
\(468\) distinct emissions constructed in
\cref{gfp:b:chap:semillas-emisor-54}, and let
\[
 \Psi:\mathcal U_6\longrightarrow\mathbb F_3^6,
 \qquad
 \Psi(U)=((-U_j)\bmod3)_{j=1}^{6}.
\]
For \(w\in\Psi(\mathcal U_6)\), we recall the invariants
\[
 n(w)=\#\Psi^{-1}(w),
 \qquad
 M(w)=\sum_{U\in\Psi^{-1}(w)}\operatorname{mult}(U),
\]
and the occupation vector
\[
 \nu(w)=
 \bigl(\#\{j:w_j=0\},\#\{j:w_j=1\},\#\{j:w_j=2\}\bigr).
\]

\begin{theorem}[Orbital selection of the self-scale]
\label{gfp:b:thm:seleccion-orbital-phi}
The dihedral action on \(\Psi(\mathcal U_6)\) contains a unique orbit
that satisfiis the self-scale predicate
\[
 n(w)=1,
 \qquad
 M(w)=144,
 \qquad
 \nu(w)=(2,2,2).
\]
The internal caliber
\[
 \operatorname{diag}_c=6,
 \qquad
 \operatorname{card}=ES
\]
select in that orbit the word
\[
 \boxed{w_6^{\varphi}=121200.}
\]
Its only emission is
\[
 \boxed{U_{\varphi}=(830,943,830,109,252,252),}
\]
of multiplicity \(144\), multiplicative signature \(555933\), additive
orientation \(ES\), multiplicative support of cardinality six and residual type
\(A\).
\end{theorem}

\begin{proof}
The direct projection gives
\[
 U_{\varphi}\bmod3=(2,1,2,1,0,0),
 \qquad
 -U_{\varphi}\bmod3=(1,2,1,2,0,0).
\]
The enumeration is carried out over the \(243\) visible words and their
\(43\) orbits, not over a selection of examples. The joint predicate
\(n=1\), \(M=144\), \(\nu=(2,2,2)\) leaves one orbit. Within it, the
declared diagonal and orientation leave one emission. At this stage neither
the equation \(x^2-x-1=0\), nor a decimal table, nor the
conventional name of the constant has been introduced.
\end{proof}

\section{Blockwise lifting and \(1^{\mathrm a}\) coinductive boundary}

Let \(b_1=w_6^{\varphi}\), and let
\(\varepsilon=(0,0,1,1)\). With the matrices \(L_0,L_1\) constructed in
\cref{gfp:b:chap:orbitas-elevacion-r36}, we define
\[
 b_{k+1}=b_kL_{\varepsilon_k}\pmod3
 \qquad(1\leq k\leq4),
\]
and concatenate
\[
 w_{6(k+1)}^{\varphi}=w_{6k}^{\varphi}\mathbin{\|}b_{k+1}.
\]
Matrix calculus produces, block by block,
\[
\begin{aligned}
 b_1&=121200,\\
 b_2&=112202,\\
 b_3&=121020,\\
 b_4&=010210,\\
 b_5&=010200,
\end{aligned}
\]
and, therefore,
\begin{equation}
 \boxed{
 w_{30}^{\varphi}
 =121200|112202|121020|010210|010200.}
 \label{gfp:b:eq:w30-phi}
\end{equation}
The dual neutrality of the boundary adds
\[
 u_5^{\varphi}=102011,
\]
so that
\begin{equation}
 \boxed{
 w_{36}^{\varphi}
 =121200|112202|121020|010210|010200|102011.}
 \label{gfp:b:eq:w36-phi}
\end{equation}

\begin{remark}[Two distinct uses of a visible word]
The word \(102011\) is also visible in the biography of \(e\). The
equality of a projection does not identify enriched states: the channel,
prefix, quotient, carry, phase, and memory belong to the type of the
object. This separation is mandatory before any
real character is constructed.
\end{remark}

\section{From the tritic calendar to the incidence multigraph}

\subsection{Primitive modal rule}

The primitive triticale calendar is
\[
 (+1,-1,0).
\]
We fix two visible modis, \(\Sigma\) and \(\Pi\), and the rule
\begin{equation}
 +1:m\longmapsto\Sigma,
 \qquad
 -1:m\longmapsto\Pi,
 \qquad
 0:m\longmapsto m.
 \label{gfp:b:eq:regla-modal-primitiva-phi}
\end{equation}
Under the cyclic closure condition, the modal orbit produced by
\eqref{gfp:b:eq:regla-modal-primitiva-phi} is
\[
 \boxed{(\Sigma,\Pi,\Pi).}
\]
The three consecutive pairs are
\[
 \Sigma\longrightarrow\Pi,
 \qquad
 \Pi\longrightarrow\Pi,
 \qquad
 \Pi\longrightarrow\Sigma.
\]

\begin{definition}[Chronological incidence matrix]
For \(i,j\in\{\Sigma,\Pi\}\), let
\[
 N_3(i,j)=
 \#\{r\in\mathbb Z/3\mathbb Z:(m_r,m_{r+1})=(i,j)\}.
\]
The subscript \(3\) counts instants of the primitive calendar; it does not denote a
matrix power.
\end{definition}

\begin{theorem}[Necessary descent of modal incidence]
\label{gfp:b:thm:descenso-incidencia-modal-phi}
In the order \((\Sigma,\Pi)\),
\begin{equation}
 \boxed{
 N_3=
 \begin{pmatrix}
 0&1\\
 1&1
 \end{pmatrix}.}
 \label{gfp:b:eq:N3-phi}
\end{equation}
By exact repetition of the calendar,
\begin{equation}
 \boxed{
 N_9=3N_3,
 \qquad
 N_{27}=9N_3,
 \qquad
 N_{108}=36N_3.}
 \label{gfp:b:eq:incidencias-9-27-108-phi}
\end{equation}
These three equalities are equalities of edge counts; they do not assert
\(N_9=N_3^3\), \(N_{27}=N_3^9\) or \(N_{108}=N_3^{36}\).
\end{theorem}

\begin{proof}
\(\Sigma\to\Sigma\) does not occur; each of the other three pairs occurs
once. This determines the four entries of \(N_3\). The calendars of
lengths \(9\), \(27\), and \(108\) contain respectively \(3\), \(9\),
and \(36\) copies of the primitive block. Each repetition multiplies all the
counts by the same integer, which proves
\eqref{gfp:b:eq:incidencias-9-27-108-phi}.
\end{proof}

\subsection{Areal-volumetric publication}

The modal label and the geometric sheet are distinct objects. The publication chart is
\[
 \mathfrak g_{\rm av}(\Sigma)=\mathscr H_{\rm ar},
 \qquad
 \mathfrak g_{\rm av}(\Pi)=\mathscr H_{\rm vol}.
\]
Transporting \eqref{gfp:b:eq:N3-phi} through \(\mathfrak g_{\rm av}\), we obtain
the incidence operator
\begin{equation}
 \boxed{
 F_{\rm av}=
 \begin{pmatrix}
 0&1\\
 1&1
 \end{pmatrix}.}
 \label{gfp:b:eq:Fav-derivada-phi}
\end{equation}
Thus, the two cross-transitions, the unique volumetric self-retention, and the
absence of areal self-retention are derived from the calendar; they are not four
independent inputs.

\begin{remark}[Scope of the descent]
Forgetting the elementary phase does not by itself produce a strong Markov
quotient of the Topological Phase Kernel, because the visible mode does not determine
which of the three internal operators will act at the next step. What
does descend without an additional hypothesis is the edge multigraph of a
complete block. Its minimal memory-one envelope has adjacency
matrix \(F_{\rm av}\).
\end{remark}

\section{Invariant positive ray and uniqueness of self-scaling}

\begin{theorem}[Perron ray derived]
\label{gfp:b:thm:rayo-perron-phi}
The positive cone of \(F_{\rm av}\) contains a unique eigenray. If we
normalize it as \(v=(1,x)^{\mathsf T}\), then
\[
 x^2-x-1=0,
 \qquad
 x>0,
\]
and, therefore,
\begin{equation}
 \boxed{x=\frac{1+\sqrt5}{2}=\varphi.}
 \label{gfp:b:eq:phi-perron}
\end{equation}
\end{theorem}

\begin{proof}
The equation \(F_{\rm av}v=\lambda v\) gives, coordinate by coordinate,
\[
 x=\lambda,
 \qquad
 1+x=\lambda x.
\]
Eliminating \(\lambda\) yields \(x^2-x-1=0\). Its roots are \((1\pm\sqrt5)/2\), and only the positive root lies in the interior of the cone. The matrix \(F_{\rm av}\) is primitive, since \(F_{\rm av}^2\) has all its entries positive; the Perron--Frobenius theorem ensures that this positive ray is simple and unique. The name \(\varphi\) is assigned only after this construction.
\end{proof}

\begin{corollary}[Scalar Self-Similarity Equation]
\[
 \varphi=1+\frac1\varphi,
 \qquad
 \varphi^2=\varphi+1,
 \qquad
 \varphi^{-1}=\varphi-1.
\]
Thi identitiis are consequencis of \eqref{gfp:b:eq:phi-perron}; they select neither
the seed nor the matrix.
\end{corollary}

\section{Publication of the self-scale \(\varphi\) through Fibonacci powers and Diophantine convergence}

Let \(F_0=0\), \(F_1=1\), and
\[
 F_{n+1}=F_n+F_{n-1}\qquad(n\geq1).
\]

\begin{theorem}[Exact Powers of Incidence]
\label{gfp:b:thm:potencias-fibonacci-phi}
For every \(n\geq1\),
\begin{equation}
 \boxed{
 F_{\rm av}^{\,n}=
 \begin{pmatrix}
 F_{n-1}&F_n\\
 F_n&F_{n+1}
 \end{pmatrix}.}
 \label{gfp:b:eq:potencias-Fav-fibonacci}
\end{equation}
\end{theorem}

\begin{proof}
For \(n=1\), the identity is the definition of \(F_{\rm av}\). If it holds for \(n\), right multiplication by \(F_{\rm av}\) gives
\[
 \begin{pmatrix}
 F_n&F_{n-1}+F_n\\
 F_{n+1}&F_n+F_{n+1}
 \end{pmatrix}
 =
 \begin{pmatrix}
 F_n&F_{n+1}\\
 F_{n+1}&F_{n+2}
 \end{pmatrix},
\]
which is the case \(n+1\).
\end{proof}

\begin{theorem}[Cassini's Identity]
\label{gfp:b:thm:cassini-phi}
For every \(n\geq1\),
\begin{equation}
 \boxed{F_{n+1}F_{n-1}-F_n^2=(-1)^n.}
 \label{gfp:b:eq:cassini-phi}
\end{equation}
\end{theorem}

\begin{proof}
The determinant of \(F_{\rm av}\) is \(-1\). Taking determinants in
\eqref{gfp:b:eq:potencias-Fav-fibonacci},
\[
 F_{n-1}F_{n+1}-F_n^2
 =\det(F_{\rm av}^{\,n})=(-1)^n.
\]
\end{proof}

We define \(r_n=F_{n+1}/F_n\) for \(n\geq1\). From Cassini it follows that
\[
 r_n-r_{n-1}
 =\frac{F_{n+1}F_{n-1}-F_n^2}{F_nF_{n-1}}
 =\frac{(-1)^n}{F_nF_{n-1}}.
\]
Therefore, the approximations alternate around the Perron ray.

\begin{proposition}[Perron Rational Gaps]
\label{gfp:b:prop:horquillas-perron-phi}
For \(m\geq1\),
\begin{equation}
 \frac{F_{2m+2}}{F_{2m+1}}
 <\varphi<
 \frac{F_{2m+1}}{F_{2m}},
 \label{gfp:b:eq:horquilla-fibonacci-phi}
\end{equation}
and the width is
\begin{equation}
 \boxed{
 \frac{F_{2m+1}}{F_{2m}}
 -\frac{F_{2m+2}}{F_{2m+1}}
 =\frac{1}{F_{2m}F_{2m+1}}.}
 \label{gfp:b:eq:anchura-fibonacci-phi}
\end{equation}
\end{proposition}

\begin{proof}
The preceding alternation places the even quotients below and the odd ones
above \(\varphi\). The formula for the width is
\eqref{gfp:b:eq:cassini-phi} with \(n=2m+1\), after putting both fractions over a
common denominator. Since \(F_n\to\infty\), the width tends to zero.
\end{proof}

\section{Memory incidence of order \(1\) and Parry measure of the symbolic envelope}

Let \(X_{\rm av}\) be the smallest subshift of finite type over
\(\{\mathscr H_{\rm ar},\mathscr H_{\rm vol}\}\) that admits exactly
the three pairs observed in the modal calendar. Its adjacency matrix is
\(F_{\rm av}\).

\begin{theorem}[Transition and Parry Measure]
\label{gfp:b:thm:parry-phi}
Let \(r=(1,\varphi)^{\mathsf T}\). The Parry stochastic matrix is
\begin{equation}
 \boxed{
 P_{ij}=
 \frac{(F_{\rm av})_{ij}r_j}{\varphi r_i}
 =
 \begin{pmatrix}
 0&1\\
 \varphi^{-2}&\varphi^{-1}
 \end{pmatrix}.}
 \label{gfp:b:eq:parry-matriz-phi}
\end{equation}
The measure estacionaria of vertices is proporcional to
\((1,\varphi^2)\), hence
\begin{equation}
 \boxed{
 \frac{\mu_{\rm ar}}{\mu_{\rm vol}}=\varphi^{-2}.}
 \label{gfp:b:eq:parry-razon-phi}
\end{equation}
\end{theorem}

\begin{proof}
The sum of row \(i\) of \(P\) is
\[
 \frac{(F_{\rm av}r)_i}{\varphi r_i}=1.
\]
Since \(F_{\rm av}\) is symmetric, the left and right eigenvectors
coincide. The stationary weights are proportional to their coordinatewise
products, that is, to \((1,\varphi^2)\). The ratio between the
two weights is \(\varphi^{-2}\).
\end{proof}

\begin{remark}[Three measures that must not be confused]
The chronological frequency of the primitive orbit is
\[
 \nu_{\rm cyc}(\mathscr H_{\rm ar})=\frac13,
 \qquad
 \nu_{\rm cyc}(\mathscr H_{\rm vol})=\frac23.
\]
The uniform measure \(\tau_{468}\) lives on the emission catalog. The
Parry measure lives on admissible histories of arbitrary length. The
first counts instants, the second counts emissions, and the third weights
dynamic cylinders. None is an approximation of the others.
\end{remark}

\section{Multiplicative character and cylindrical conservation}

For an admissible route \(\gamma:i\to j\) of length \(n\), we define
\begin{equation}
 \chi_P(\gamma)=\varphi^n\frac{r_i}{r_j}.
 \label{gfp:b:eq:caracter-parry-phi}
\end{equation}

\begin{proposition}[Concatenation multiplicativity]
\label{gfp:b:prop:multiplicatividad-parry-phi}
If \(\gamma_1:i\to j\) and \(\gamma_2:j\to k\), then
\[
 \chi_P(\gamma_2\circ\gamma_1)
 =\chi_P(\gamma_2)\chi_P(\gamma_1).
\]
On to closed route of length \(n\),
\[
 \chi_P(\gamma)=\varphi^n.
\]
\end{proposition}

\begin{proof}
Length adds and the intermediate coordinate telescopes:
\[
 \varphi^{n_1+n_2}\frac{r_i}{r_k}
 =\left(\varphi^{n_1}\frac{r_i}{r_j}\right)
  \left(\varphi^{n_2}\frac{r_j}{r_k}\right).
\]
If \(i=k\), the quotient terminal is one.
\end{proof}

We normalize the left vector as
\[
 \ell=\frac{r}{1+\varphi^2},
 \qquad
 \ell^{\mathsf T}r=1.
\]
The measure of a cylinder is
\[
 \mu[x_0\cdots x_n]
 =\ell_{x_0}r_{x_n}\varphi^{-n}.
\]
Therefore,
\begin{equation}
 V_n(x)
 :=\frac{\mu[x_0]}{\mu[x_0\cdots x_n]}
 =\varphi^n\frac{r_{x_0}}{r_{x_n}}
 =\chi_P(x_0\cdots x_n).
 \label{gfp:b:eq:volumen-cilindrico-phi}
\end{equation}
For to typed dimension \(D>0\), the scale
\[
 a_n^{(D)}(x)=V_n(x)^{1/D}
\]
satisfies the exact law
\begin{equation}
 \boxed{
 \bigl(a_n^{(D)}(x)\bigr)^D
 \mu[x_0\cdots x_n]=\mu[x_0].}
 \label{gfp:b:eq:conservacion-cilindrica-phi}
\end{equation}

In three dimensions and over closed returns,
\[
 \frac{a_9}{a_0}=\varphi^3,
 \qquad
 \frac{a_{27}}{a_0}=\varphi^9,
 \qquad
 \frac{a_{108}}{a_0}=\varphi^{36}.
\]
If \(A_k=a_{9k}\), then
\begin{equation}
 A_{k+1}-2A_k+A_{k-1}
 =A_{k-1}(\varphi^3-1)^2>0.
 \label{gfp:b:eq:convexidad-retornos-phi}
\end{equation}
This is convexity in return depth; it is not interpreted as a
kinematic quantity without an additional duration chart.

\section{Quadratic memory and contracted branch}

The eigenvalues of \(F_{\rm av}\) are
\[
 \varphi,
 \qquad
 -\varphi^{-1}.
\]
The contracted branch changes orientation at each iteration. To preserve
that sign before squaring, we pass to
\(\operatorname{Sym}^2(\mathbb R^2)\). In the basis
\((x^2,2xy,y^2)\),
\begin{equation}
 \operatorname{Sym}^2(F_{\rm av})=
 \begin{pmatrix}
 0&0&1\\
 0&1&2\\
 1&1&1
 \end{pmatrix}.
 \label{gfp:b:eq:sym2-Fav-phi}
\end{equation}

\begin{theorem}[Spectrum of the second-order memory]
\label{gfp:b:thm:espectro-sym2-phi}
\[
 \det\bigl(tI-\operatorname{Sym}^2(F_{\rm av})\bigr)
 =(t+1)(t^2-3t+1),
\]
and
\begin{equation}
 \boxed{
 \operatorname{Spec}\bigl(\operatorname{Sym}^2(F_{\rm av})\bigr)
 =\{\varphi^2,-1,\varphi^{-2}\}.}
 \label{gfp:b:eq:espectro-sym2-phi}
\end{equation}
The only positive stable mode is \(\varphi^{-2}\).
\end{theorem}

\begin{proof}
The eigenvalues of a degree-two symmetric power are the products
\(\lambda_1^2\), \(\lambda_1\lambda_2\), and \(\lambda_2^2\). With
\(\lambda_1=\varphi\) and \(\lambda_2=-\varphi^{-1}\), one obtains
\(\varphi^2\), \(-1\), and \(\varphi^{-2}\). The first exceeds one, the
second has modulus one, and the third belongs to \((0,1)\).
\end{proof}

On a closed route with \(V_n=\varphi^n\), the stable mode satisfies
\begin{equation}
 \boxed{M_n=\varphi^{-2n}=V_n^{-2}.}
 \label{gfp:b:eq:modo-estable-phi}
\end{equation}

\section{Normalized hyperbolic realization of the self-scale: \(\exp(2\log\varphi\,J_h)=F_{\mathrm{av}}^2\)}

Let
\[
 A=F_{\rm av}^{\,2}=
 \begin{pmatrix}1&1\\1&2\end{pmatrix},
 \qquad
 J_h=\frac{2A-3I}{\sqrt5}.
\]
Then
\[
 J_h^2=I,
 \qquad
 A=\frac32I+\frac{\sqrt5}{2}J_h.
\]
Since
\[
 \cosh(2\log\varphi)=\frac32,
 \qquad
 \sinh(2\log\varphi)=\frac{\sqrt5}{2},
\]
we obtain
\begin{equation}
 \boxed{
 \exp(2\log\varphi\,J_h)=F_{\rm av}^{\,2}.}
 \label{gfp:b:eq:realizacion-hiperbolica-phi}
\end{equation}
The self-scaling constant is therefore simultaneously the Perron
eigenvalue of the incidence and the exponential parameter of its normalized
hyperbolic realization.

\section{Characterization, brackets and prolongation}

The finite construction determines the positive ray of \(F_{\rm av}\). The
subsequent analytic characterization is
\[
 x^2-x-1=0,
 \qquad x>0.
\]
The brackets \eqref{gfp:b:eq:horquilla-fibonacci-phi} have rational endpoints
and explicit width \eqref{gfp:b:eq:anchura-fibonacci-phi}. For each
nonadic level \(K\), we choo the least \(m=m_{\varphi}(K)\) who bracket remains
contained in a single subcylinder among the \(729\) elements of the
ambient fiber. Minimality converts localization into a function, not into to
tacit choice.

The selected cylinders are compatible with the truncations, and their
diameters tend to zero. Their intersection contains a unique real number; the
\cref{gfp:b:thm:rayo-perron-phi} identifies it with \(\varphi\). The Fibonacci
quotients are therefore rational certificates of localization, not decimal data
used to decide the seed.

\Needspace{22\baselineskip}
\section{Internal certificate of uniqueness, compatibility and convergence of Perron self-scaling}

\begin{remark}[Certified and validity control: Refutation of the self-biography]
The construction is refuted if any of the following events occurs:
\begin{enumerate}
 \setlength{\itemsep}{0.15\baselineskip}
 \item the finite predicate selects more than one orbit or more than one emission;
 \item the recurrence \(L_0,L_0,L_1,L_1\) does not produce
       \eqref{gfp:b:eq:w30-phi};
 \item the calendar does not produce exactly the three pairs used in
       \eqref{gfp:b:eq:N3-phi};
 \item a count multiple \(N_{9},N_{27},N_{108}\) is confused with a
       power of \(N_3\);
 \item \(F_{\rm av}\) has another positive eigenray;
 \item \eqref{gfp:b:eq:potencias-Fav-fibonacci} or Cassini’s identity fails;
 \item the rational brackets are not nested or do not contract their width;
 \item the chronological frequency, the uniform measure on the
       catalog or the Parry measure are identified;
 \item a decimal table is involved before the Perron ray is constructed;
 \item a selected cylinder does not truncate to the previous cylinder.
\end{enumerate}
\end{remark}
